\documentclass[12pt,oneside]{book}

% ============================================================================
% METADATOS Y CONFIGURACIÓN INICIAL
% ============================================================================

\title{Apuntes de Derecho Procesal Civil}
\author{Isaac Edgar Camacho Ocampo}
\date{2024}

% ============================================================================
% PAQUETES ESENCIALES
% ============================================================================

\usepackage[utf8]{inputenc}
\usepackage[T1]{fontenc}
\usepackage[spanish,es-nodecimaldot]{babel}

\usepackage[
    top=2.5cm,
    bottom=2.5cm,
    left=3cm,
    right=2.5cm,
    headheight=15pt
]{geometry}

\usepackage{amsmath}
\usepackage{amssymb}
\usepackage{mathtools}

\usepackage{graphicx}
\usepackage{float}
\usepackage{caption}
\usepackage{subcaption}

\usepackage{booktabs}
\usepackage{array}
\usepackage{longtable}
\usepackage{tabularx}

\usepackage{enumitem}

\usepackage{xcolor}
\usepackage[most]{tcolorbox}

\usepackage{fancyhdr}

\usepackage{ragged2e}

\usepackage[
    colorlinks=true,
    linkcolor=blue!50!black,
    citecolor=green!40!black,
    urlcolor=blue!60!black,
    pdftitle={Apuntes de Derecho Procesal Civil},
    pdfauthor={Isaac Edgar Camacho Ocampo},
    pdfsubject={Apuntes universitarios},
    bookmarks=true
]{hyperref}

% ============================================================================
% VARIABLES DE METADATOS
% ============================================================================

\newcommand{\universidad}{Universidad de Buenos Aires}
\newcommand{\facultad}{Facultad de Derecho}
\newcommand{\materia}{Derecho Procesal Civil}
\newcommand{\titulo}{Unidad 18: Resoluciones Judiciales y Cosa Juzgada}
\newcommand{\autor}{Isaac Edgar Camacho Ocampo}
\newcommand{\anio}{2024}

% ============================================================================
% CONFIGURACIÓN DE MÁRGENES Y ESPACIADO
% ============================================================================

\setlength{\parindent}{0pt}
\setlength{\parskip}{0.5em}

\setlist[itemize]{
    topsep=0.3em,
    itemsep=0.2em
}

\setlist[enumerate]{
    topsep=0.3em,
    itemsep=0.2em
}

\clubpenalty=10000
\widowpenalty=10000

% ============================================================================
% CONFIGURACIÓN DE ENCABEZADOS Y PIES
% ============================================================================

\pagestyle{fancy}
\fancyhf{}

\fancyhead[L]{\nouppercase{\leftmark}}
\fancyhead[R]{\materia}

\fancyfoot[C]{\thepage}

\renewcommand{\headrulewidth}{0.4pt}
\renewcommand{\footrulewidth}{0pt}

\fancypagestyle{plain}{
    \fancyhf{}
    \fancyfoot[C]{\thepage}
    \renewcommand{\headrulewidth}{0pt}
}

% ============================================================================
% CONFIGURACIÓN DE ÍNDICE
% ============================================================================

\setcounter{tocdepth}{2}

% ============================================================================
% COMANDOS MATEMÁTICOS Y UTILIDADES
% ============================================================================

\newcommand{\abs}[1]{\left|#1\right|}
\newcommand{\norm}[1]{\left\lVert#1\right\rVert}
\newcommand{\vect}[1]{\mathbf{#1}}
\newcommand{\deriv}[2]{\frac{d#1}{d#2}}
\newcommand{\pderiv}[2]{\frac{\partial#1}{\partial#2}}

% ============================================================================
% DEFINICIÓN DE CAJAS ACADÉMICAS
% ============================================================================

\newtcolorbox{definition}[1]{
    enhanced,
    breakable,
    title={Definición: #1},
    fonttitle=\bfseries,
    colback=blue!3,
    colframe=blue!50!black,
    boxrule=0.8pt,
    arc=2mm
}

\newtcolorbox{important}{
    enhanced,
    breakable,
    title={Importante},
    fonttitle=\bfseries,
    colback=yellow!8,
    colframe=orange!70!black,
    boxrule=0.8pt,
    arc=2mm
}

\newtcolorbox{example}[1]{
    enhanced,
    breakable,
    title={Ejemplo: #1},
    fonttitle=\bfseries,
    colback=green!4,
    colframe=green!40!black,
    boxrule=0.8pt,
    arc=2mm
}

\newtcolorbox{formula}[1]{
    enhanced,
    breakable,
    title={Fórmula / Regla: #1},
    fonttitle=\bfseries,
    colback=purple!4,
    colframe=purple!50!black,
    boxrule=0.8pt,
    arc=2mm
}

\newtcolorbox{exercise}[1]{
    enhanced,
    breakable,
    title={Ejercicio: #1},
    fonttitle=\bfseries,
    colback=gray!5,
    colframe=gray!60!black,
    boxrule=0.8pt,
    arc=2mm
}

\newtcolorbox{note}{
    enhanced,
    breakable,
    title={Nota},
    fonttitle=\bfseries,
    colback=white,
    colframe=gray!50,
    boxrule=0.6pt,
    arc=2mm
}

% ============================================================================
% RUTAS DE IMÁGENES
% ============================================================================

\graphicspath{
    {./img/}
    {./graficos/}
    {./figuras/}
}

% ============================================================================
% INICIO DEL DOCUMENTO
% ============================================================================

\begin{document}

% ============================================================================
% PORTADA
% ============================================================================

\begin{titlepage}

\thispagestyle{empty}

\begin{center}

\vspace*{1cm}

{\large \textsc{\universidad}}

\vspace{0.2cm}

{\large \textsc{\facultad}}

\vspace{3cm}

{\Huge \textbf{\materia}}

\vspace{1cm}

{\LARGE \titulo}

\vspace{0.8cm}

\textit{El dominio de los mecanismos procesales es la base para ejercer adecuadamente el derecho de defensa en todos los órdenes.}

\vspace{1.2cm}

\rule{0.70\textwidth}{0.5pt}

\vspace{0.5cm}

{\large Apuntes de estudio}

\vspace{0.5cm}

\rule{0.70\textwidth}{0.5pt}

\vfill

{\large \autor}

\vspace{0.3cm}

{\large \anio}

\end{center}

\vfill

\begin{flushleft}
\textit{Este es un modesto aporte a los estudiantes de ``Derecho Procesal Civil'' no pretende sustituir el material oficial de la materia.}
\end{flushleft}

\end{titlepage}

% ============================================================================
% ÍNDICE
% ============================================================================

\tableofcontents

\newpage

% ============================================================================
% CAPÍTULO 1: RESOLUCIONES JUDICIALES
% ============================================================================

\chapter{Resoluciones Judiciales}

\section{Concepto y Clasificación}

Las resoluciones judiciales constituyen las diversas formas en que el tribunal se pronuncia ante las cuestiones que le son planteadas. El tribunal debe pronunciarse desde actos de gran complejidad —como la sentencia definitiva— hasta actos procesales muy sencillos, como las providencias simples.

El Código Procesal agrupa las resoluciones judiciales de acuerdo a su complejidad y al objetivo que tiene cada una. Todas las resoluciones judiciales son actos procesales con características definidas: deben ser redactadas en idioma español, en tinta negra, fechadas y firmadas por la autoridad competente del tribunal.

La clasificación general es la siguiente:

\begin{itemize}
    \item \textbf{Providencias simples}: actos procesales de mínima complejidad que no requieren motivación.
    \item \textbf{Resoluciones interlocutorias}: decisiones sobre cuestiones incidentales que requieren motivación y sustanciación.
    \item \textbf{Sentencias homologatorias}: resoluciones que homologan determinados modos anormales de terminación del proceso.
    \item \textbf{Sentencias definitivas}: los actos procesales de mayor complejidad que ponen fin al proceso de manera normal.
\end{itemize}

\subsection{La Importancia de la Clasificación}

Conocer ante qué tipo de resolución se está es fundamental para el litigante, porque cada resolución judicial tiene determinados remedios o recursos que pueden ser interpuestos contra ella en un plazo específico.

\begin{important}
No puede utilizarse cualquier recurso contra cualquier resolución. Cada acto procesal tiene sus propios remedios disponibles. Utilizar un recurso inadecuado implicaría perder la posibilidad de corrección.
\end{important}

\begin{example}{La analogía del cáñón y el mosquito}
Uno no puede matar un mosquito con un cañón, ni con una pistola ni con una ametralladora. Con los recursos procesales sucede exactamente lo mismo: no puede atacarse ciertos actos del proceso con un recurso de mayor envergadura cuando a veces se necesita únicamente una aclaratoria o una revocatoria. La revocatoria no se puede interponer contra cualquier tipo de resolución judicial. Para saber qué recursos están disponibles, primero debe saberse qué tipo de acto se tiene delante.
\end{example}

\section{Providencias Simples}

Las providencias simples son resoluciones que hacen al normal desarrollo y avance del proceso, que no agregan nada relevante sino únicamente hacen que el proceso siga avanzando en su trámite.

Ejemplos de providencias simples incluyen:

\begin{itemize}
    \item Agregar algún documento o prueba que ha sido ordenada.
    \item Tener por contestada la demanda.
    \item Correr una vista o un traslado.
    \item Ordenar una notificación.
    \item Autorizar sacar fotocopias o consultar la causa.
\end{itemize}

\subsection{Plazo de Dictado y Autoridad de Firma}

El Código Procesal establece en el artículo 34 los plazos para el dictado de providencias simples. El plazo del tribunal es de \textbf{tres días} para expedirse en aquellas cuestiones que hacen al mero trámite de la causa.

El Código autoriza al secretario del tribunal, al pro-secretario e incluso al jefe de despacho a firmar determinadas providencias simples para disminuir la carga de trabajo del juez. En general, pueden firmar aquellas vinculadas a:

\begin{itemize}
    \item Agregación o devolución de escritos.
    \item Devolución de escritos presentados fuera de término.
    \item Agregación de pruebas y documentos.
    \item Vistas que deben darse al ministerio público.
\end{itemize}

\subsection{Providencias que Requieren Firma del Juez}

Sin embargo, determinadas providencias simples sí o sí deben ser suscritas por el juez:

\begin{itemize}
    \item El reconocimiento de la condición de parte.
    \item Tener por contestada la demanda.
    \item Aquellas que tienen una relevancia mayor en el proceso.
\end{itemize}

\begin{example}{El daño de una providencia simple mal resuelta}
Si la providencia simple ``téngase por contestada la demanda'' es rechazada por considerarla extemporánea, y en lugar de ello el juez firma que la contestación fue presentada fuera de término, existe un daño irreparable al litigante. Por eso es tan importante saber que se está ante una providencia simple, porque el remedio disponible es la revocatoria, que a veces es el único que puede solucionar ese problema.
\end{example}

\section{Resoluciones Interlocutorias}

Las resoluciones interlocutorias son decisiones que el tribunal dicta en medio del proceso sobre cuestiones incidentales. Quizás lo más difícil para estudiantes y abogados es diferenciar una providencia simple de una resolución interlocutoria.

\subsection{El Elemento Diferenciador: la Sustanciación}

El elemento diferenciador clave es la \textbf{sustanciación}. Esta palabra no se refiere al traslado (uso forense coloquial), sino a la sustancia, al contenido que tiene esa resolución. Es decir, la motivación y las razones que el juez debe desarrollar para tomar su decisión.

Las resoluciones interlocutorias están decidiendo algo: pueden hacer lugar o rechazar una excepción, convocar o rechazar la convocatoria de un tercero, ordenar o rechazar una medida cautelar. Para llegar a esa decisión, el juez tiene que motivar y dar razones de por qué llega a ella.

\begin{note}
La palabra `interlocutorio' proviene del latín `inter', que significa `en medio del debate'. Son decisiones que se toman en medio de la discusión principal.
\end{note}

\subsection{Ejemplo Demostrativo}

\begin{example}{La medida cautelar confirma que sustanciación no es traslado}
Cuando una parte pide una medida cautelar, no hay traslado previo. Sin embargo, la decisión que toma el juez no es una providencia simple: es una sentencia interlocutoria. Esto demuestra que el código no se refiere a la sustanciación como traslado, sino como sustancia, como materia, como motivación del acto. De otro modo, uno debería concluir que una medida cautelar es una providencia simple, lo cual es incorrecto.
\end{example}

\subsection{Plazos para Dictar Resoluciones Interlocutorias}

Los plazos para dictar estas resoluciones son:

\begin{itemize}
    \item \textbf{Primera instancia}: 10 días.
    \item \textbf{Segunda instancia}: 15 días.
\end{itemize}

\subsection{Capacidad Resolutiva}

Hay resoluciones interlocutorias que pueden llegar a poner fin al proceso —como una sentencia que hace lugar a una excepción de prescripción o a una caducidad de la instancia— aunque lo hacen de un modo anormal, no del modo normal que solo se alcanza mediante la sentencia definitiva.

\section{Sentencias Homologatorias}

Las sentencias homologatorias son resoluciones vinculadas a determinados mecanismos que constituyen modos anormales de terminación del proceso: un allanamiento, un acuerdo conciliatorio, un desistimiento.

El Código prevé que se dicte una resolución que no llega a ser una sentencia definitiva, pero que tampoco es una providencia simple. Tienen menos complejidad en su sustancia pero no dejan de tenerla, porque el Código le está exigiendo al juez que evalúe si el acto de la parte es correcto y si tiene las condiciones necesarias para lograr el desistimiento, el allanamiento o el acuerdo conciliatorio.

\begin{important}
Los acuerdos transaccionales homologados tienen efectos de cosa juzgada.
\end{important}

\section{Sentencia Definitiva de Primera Instancia}

La sentencia definitiva es el acto más complejo que el Código prevé dentro de las resoluciones judiciales. En el artículo 163 del Código Procesal Civil y Comercial se encuentran los nueve incisos que describen las distintas condiciones que debe tener una sentencia definitiva.

\subsection{Estructura Formal}

Tanto la resolución interlocutoria como la sentencia definitiva suelen comenzar con:

\begin{enumerate}
    \item La fecha y la palabra ``Autos'' o ``Autos y Vistos''.
    \item El resumen de los planteos de las partes (resultando).
    \item El ``Y considerando'' (donde el tribunal evalúa los planteos).
    \item La parte resolutiva (el fallo propiamente dicho).
\end{enumerate}

\subsection{El Deber de Motivación}

\begin{important}
La motivación de la decisión judicial no es solo explicar por qué se llega a una resolución determinada. La motivación de las decisiones judiciales sirve para que las partes y —de modo muy indirecto— el resto de la población controlen los actos de gobierno. Las sentencias y las decisiones judiciales son actos de gobierno: el Poder Judicial es uno de los poderes del Estado.
\end{important}

Además, la motivación es la base del control recursivo: para interponer una apelación contra una sentencia, el Código exige una crítica concreta y razonada de la decisión que se quiere modificar. Esa crítica solo puede construirse sobre los considerandos, que es la parte donde el juez brinda las explicaciones de por qué llegó al resultado.

\subsection{Presunciones e Indiciaria}

Según el artículo 163 del Código Procesal, la presunción no es un medio de prueba, sino una forma de dirimir la cuestión cuando el hecho principal no está probado, pero se tiene una serie de hechos que, por su importancia, gravedad y concordancia, hacen que se deba presumir que aquel hecho que no está probado de forma directa ha sucedido. Es, en definitiva, un mecanismo de razonamiento que la ley brinda al juez para poder dictar sentencia.

El juez civil no puede mantener silencio y no dictar sentencia porque un hecho no está probado; puede echar mano a los razonamientos de la prueba presuncional o indiciaria.

\subsection{Conducta Procesal de las Partes}

\begin{example}{Conducta procesal como elemento de valoración}
Si una parte está en condiciones de traer determinado medio de prueba y omite hacerlo sin explicar por qué, lo más razonable es pensar que ese medio de prueba le era desfavorable. No puede premiarse una conducta que obstaculiza la finalidad del proceso, que es averiguar cómo sucedieron los hechos para determinar si hay que condenar o absolver al demandado.
\end{example}

\subsection{Congruencia y Límites de la Decisión}

El inciso sexto del artículo 163 establece la decisión expresa, positiva y precisa de conformidad con las pretensiones deducidas en el juicio. Esto no es más ni menos que el principio dispositivo y la congruencia.

\begin{important}
El juez cuando dicta la sentencia tiene que hacerse cargo de aquellas cuestiones que han sido introducidas. Pero esto también le marca un límite: el juez no puede dar más de lo que se le pidió, ni algo distinto de lo que se le pidió. Hacerlo implicaría violar la congruencia y el derecho de defensa de la parte demandada.
\end{important}

\begin{example}{La incapacidad no reclamada}
En un proceso donde no se reclama el resarcimiento de determinada incapacidad, el juez no podría concederla porque le parezca razonable. Si lo hiciera, estaría violando la congruencia y dando por más allá de lo que ha sido solicitado por los litigantes, además de violar el derecho de defensa de la parte demandada que no ha podido denegar que tal incapacidad existió.
\end{example}

\subsection{Hechos Sobrevinientes}

Respecto de los hechos sobrevinientes, en juicios que duran dos, tres o cuatro años pueden suceder hechos que modifiquen la relación jurídica que es objeto del proceso, la extingan, o que hagan que la realidad exterior exija que la decisión del juez se adecúe a ese cambio.

\begin{example}{El cuadro de Picasso destruido}
Si en un proceso se reclama la restitución de un cuadro pintado por Picasso, y durante el juicio ese cuadro deja de existir (porque se destruyó o desapareció), el juez no deberá condenar a su restitución en la sentencia (orden que no podría cumplirse), sino que deberá reemplazarla por una suma de dinero que compense el valor del bien que debía serle restituido a esa parte.
\end{example}

\subsection{Las Costas}

El inciso octavo del artículo 163 establece que el juez debe, tanto en la sentencia definitiva como en las interlocutorias, establecer lo que son las costas y los honorarios.

\begin{definition}{Costas}
Las costas comprenden los gastos que ha incurrido la parte para llevar adelante ese reclamo. No solo los honorarios del abogado, sino todo otro gasto en el cual haya debido incurrir (sellados, obtención de documentos, etc.).
\end{definition}

\subsection{Regla del Vencido}

La regla que establece el Código para la condena en costas es la del vencido: es decir, si hay una parte que es vencida en el proceso, es la que debe hacerse cargo del pago de las costas. Si la demanda prospera, en general el vencido es el demandado; si la demanda es rechazada, el vencido será el actor.

Sin embargo, como toda regla, también existen excepciones. Una posibilidad es la distribución en el orden causado (cada parte paga sus propios gastos; los comunes, por mitades).

\begin{important}
Debe tenerse cuidado con los argumentos para eximirse de costas. Si uno dijera ``se creyó con derecho a litigar'' como fundamento para eximirse de costas, ese argumento es endeble. Si no se hubiera creído con derecho a litigar, ya estaría incurriendo en la sanción por temeridad y malicia. Es bastante excepcional esta posibilidad, especialmente en los casos en que se discuten derechos patrimoniales. En cambio, en cuestiones de familia —como en los alimentos— es muy usual ver la regla de la distribución en el orden causado para no afectar el monto del alimento que se está reclamando.
\end{important}

% ============================================================================
% CAPÍTULO 2: SEGUNDA INSTANCIA Y ACTUACIÓN POSTERIOR A LA SENTENCIA
% ============================================================================

\chapter{Segunda Instancia y Actuación Posterior a la Sentencia}

\section{Sentencia Definitiva de Segunda Instancia}

En la cámara —segunda instancia— hay diferencias en el modo en que tienen que ser elaboradas las resoluciones. Todo lo que serían recursos que den lugar a resoluciones interlocutorias, que serán firmadas por tres vocales en cada una de las salas, van a tener que estar redactadas de modo impersonal (es decir, en persona gramatical impersonal), aunque con estructura similar a las resoluciones interlocutorias.

La sentencia definitiva de segunda instancia, en cambio, es un acto complejo que estará elaborado por un acuerdo de tres jueces: habrá un vocal preopinante y dos vocales que podrán adherir al voto del primer vocal o marcar sus diferencias (hacer disidencias). El fallo puede salir dos a uno, y no siempre el vocal preopinante tendrá el voto que determine el sentido de la decisión.

\subsection{El Acuerdo de Cámara}

El acuerdo es, en teoría, una reunión de los tres jueces para discutir las decisiones y hacer un intercambio de opiniones. En la práctica, sin embargo, no se lleva a cabo en la mayoría de los casos; lo que se hace en las cámaras es trabajar con la circulación de los expedientes de modo escrito.

\begin{important}
No puede dictarse una sentencia de segunda instancia si uno de los jueces no está presente. Podrán verse dos firmas siempre que el tercer juez esté en uso de licencia concedida o esté vacante la vocalía, pero las cámaras no tienen la misma facultad que la Corte Federal para dictar decisiones cuando alcanza la mayoría aunque los otros jueces no las suscriban.
\end{important}

\subsection{Cosa Juzgada y Modificaciones Parciales}

Una vez dictada la sentencia —sea de primera o de segunda instancia— el Código establece que esto concluye la competencia del juez o del tribunal para expedirse sobre el caso. Ya el juez no puede modificar el fallo más allá de que las partes estén o no notificadas.

Cuando una sentencia de segunda instancia modifica parcialmente la decisión de primera instancia, la cosa juzgada va a emanar de la sumatoria de los dos fallos: aquellas cuestiones que quedaron firmes por no haber sido apeladas en la sentencia de primera instancia se sumarán a las que emanan de la sentencia firme de segunda instancia.

\section{Actuación del Juez con Posterioridad a la Sentencia}

\subsection{Artículo 166 del Código Procesal}

El artículo 166 del Código Procesal establece cuál puede ser la actuación posterior del juez una vez dictada la sentencia. Básicamente, se prevé la posibilidad del juez de enmendar errores puramente numéricos o errores evidentes (errores materiales), pero en ningún caso el tribunal o el juez puede modificar lo sustancial de las decisiones ya tomadas.

\subsection{Silencio sobre Costas}

Un aspecto relevante es qué sucede cuando el tribunal omite pronunciarse sobre las costas. Durante mucho tiempo se decía que el silencio sobre costas en una sentencia implicaba que habían sido fijadas en el orden causado. Esto fue aclarado en un fallo de la Corte Federal, que dijo que el silencio de ninguna manera podría otorgarle un sentido. Lo que tiene que hacer el tribunal es expedirse sobre un tema sobre el que aún no se había expedido: la condena en costas.

\section{El Recurso de Aclaratoria}

En el segundo párrafo del artículo 166, el Código establece que dentro de los tres días y sin sustanciación se pueden corregir errores materiales y aclarar conceptos oscuros sin alterar lo sustancial.

Esto da lugar, cuando una de las partes lo solicite, al que se denomina \textbf{recurso de aclaratoria}.

\begin{note}
La aclaratoria solo puede corregir errores materiales y aclarar conceptos oscuros. No puede modificar lo sustancial de la decisión.
\end{note}

\subsection{Plazos para Solicitar Aclaratoria}

El plazo para solicitarla es de:

\begin{itemize}
    \item \textbf{Tres días} en primera instancia.
    \item \textbf{Cinco días} en segunda instancia.
\end{itemize}

% ============================================================================
% CAPÍTULO 3: COSA JUZGADA
% ============================================================================

\chapter{Cosa Juzgada}

\section{Concepto e Inmutabilidad}

El concepto de cosa juzgada implica que, contra una decisión ya sea de un tribunal de primera o de segunda instancia, no hay más recursos disponibles para interponer, o bien el plazo para apelarla ha transcurrido sin que las partes hayan deducido ninguno. En esas condiciones, puede afirmarse que la decisión adquiere la autoridad de cosa juzgada.

\begin{important}
La cosa juzgada significa que dentro del patrimonio de la persona ingresa el activo que es el derecho a hacer cumplir lo que emana de esa decisión judicial. La realidad previa deja de existir, y lo único que queda es lo que emana de la sentencia.
\end{important}

El abogado debe controlar muy celosamente lo que surge de la decisión, especialmente lo que el tribunal escribe o asienta después de la palabra ``fallo'', porque eso es lo que en el futuro se podrá exigir que se cumpla con autoridad de cosa juzgada.

\section{Cosa Juzgada Formal versus Cosa Juzgada Material}

La doctrina diferencia lo que se denomina cosa juzgada formal de cosa juzgada material:

\begin{definition}{Cosa Juzgada Material}
Cosa juzgada material es la que impide reabrir cualquier cuestión del fondo del asunto en cualquier proceso posterior. Es característica de los procesos ordinarios donde hay un conocimiento pleno del tribunal.
\end{definition}

\begin{definition}{Cosa Juzgada Formal}
Cosa juzgada formal es aquella que solo impide la discusión en el mismo proceso, pero permite iniciar un juicio ordinario posterior sobre cuestiones que no pudieron debatirse en el proceso anterior.
\end{definition}

\subsection{Ejemplo de Cosa Juzgada Formal}

\begin{example}{El juicio ejecutivo}
En un juicio ejecutivo hay cosa juzgada material si se discutió el pago, porque el pago es una excepción que se podía oponer en ese proceso. En cambio, no habrá cosa juzgada material en relación a, por ejemplo, algún vicio de la voluntad que se pudiera haber tenido cuando se suscribió ese título ejecutivo, porque esa cuestión no se podía discutir en el juicio ejecutivo. En lo que refiere a ese aspecto hay una cosa juzgada formal, y esas cuestiones cuya discusión no puede darse en el juicio ejecutivo se podrán dar en el juicio ordinario posterior.
\end{example}

\section{Límites Subjetivos de la Cosa Juzgada}

Los límites subjetivos son claros: la cosa juzgada abarcará a las personas que pudieron discutir en el marco de ese proceso previo (actor, demandado, terceros intervinientes, terceros legitimados sobre los cuales se ha discutido una relación sustancial). Pero no puede afectar a sujetos que no han participado de ese proceso, porque no han tenido la posibilidad de defenderse, ofrecer prueba y contradecir.

Esto se vincula con el tema del litisconsorcio:

\begin{important}
En tanto se trata de un litisconsorcio facultativo, nada impide que el juicio llevado adelante produzca una cosa juzgada material que afecte solo a las partes que han participado, sin que esto tenga ninguna relevancia respecto a otros sujetos que podrían haber sido litisconsortes facultativos y que no fueron citados.
\end{important}

\section{Límites Objetivos de la Cosa Juzgada}

Los límites objetivos son más difusos. Nadie tiene dudas sobre que la cuestión debatida en el primer proceso es la que ha sido juzgada. Pero cuando se lee el inciso 6º del artículo 347, se advierte que la cuestión es mucho más expansiva de lo que una primera aproximación al tema podría indicar.

Según el artículo 347, inciso 6º del Código Procesal, para que sea procedente la excepción de cosa juzgada, el examen integral de las dos contiendas debe demostrar que se trata del mismo asunto sometido a decisión judicial. Pero para que haya cosa juzgada debe ser el mismo asunto, o que por existir \textbf{continencia, conexidad, accesoriedad o subsidiariedad}, la sentencia firme ya haya resuelto lo que constituye la materia o la pretensión deducida en el nuevo juicio que se promueve.

\begin{important}
Los límites objetivos de la cosa juzgada no solo hacen a lo que fue debatido en el primer proceso, sino también a cuestiones de continencia, conexidad, accesoriedad y subsidiariedad que podrían haber sido planteadas en el primer proceso y no lo fueron. Es importante conocer esto porque la cosa juzgada puede ser declarada de oficio por el juez. Puede suceder que haya un juez que no la advierta y un demandado que tampoco oponga la excepción. En el proceso le va mejor al que más sabe y tiene más herramientas para plantear la discusión de estos supuestos.
\end{important}

\subsection{La Cosa Juzgada Implícita}

\begin{example}{El accidente de autos y el lucro cesante}
Si en un juicio por un accidente de autos se reclama solo la reparación del automóvil y se resulta ganador, ¿podría iniciarse un segundo proceso reclamando el lucro cesante (porque el auto era utilizado como taxi o remís)? Leyendo el artículo 347, inciso 6º, se podría llegar a concluir que existió lo que podría denominarse ``cosa juzgada implícita'': que el lucro cesante era un rubro que se podía pedir en el primer proceso y no se hizo, por lo que precluyó también la posibilidad de hacerlo con posterioridad, al ser una cuestión accesoria o conexa con la discutida en el primer litigio.
\end{example}

% ============================================================================
% CAPÍTULO 4: RELACIÓN ENTRE SENTENCIA CIVIL Y SENTENCIA PENAL
% ============================================================================

\chapter{Relación entre Sentencia Civil y Sentencia Penal}

\section{Prelación de la Acción Penal sobre la Civil}

Vinculado con el tema de la cosa juzgada está lo que usualmente se denomina la incidencia de la acción penal sobre la acción civil. Esto puede encontrarse en el Código Civil y Comercial a partir del artículo 1774.

\subsection{Razón de Ser de la Coordinación}

La razón de ser de esta coordinación es evitar lo que en general se encuentra en los libros como el ``escándalo jurídico'', que en definitiva es un escándalo lógico:

\begin{important}
El ``escándalo jurídico'' (o lógico) se produce cuando una sentencia en un fuero dice que el autor del hecho fue una persona, y otra sentencia dice que esa persona no fue la autora. La coordinación de los dos ordenamientos busca evitar precisamente esta colisión de decisiones.
\end{important}

\subsection{Régimen de Suspensión}

Si la acción penal precede a la acción civil o es intentada durante su curso, el dictado de la sentencia definitiva debe suspenderse en el proceso civil hasta que haya conclusión del proceso penal.

\subsection{Excepciones a la Suspensión}

Sin embargo, existen excepciones a la suspensión:

\begin{itemize}
    \item Si median causas de extinción de la acción penal.
    \item Si la dilación del procedimiento penal provoca en los hechos una frustración efectiva del derecho a ser indemnizado en el juicio civil.
    \item Si la acción civil por reparación del daño está fundada en un factor objetivo de responsabilidad. Esta es una innovación importante del Código Civil y Comercial: en estos casos nada impide que la sentencia civil se dicte aun cuando no se haya dictado sentencia penal.
\end{itemize}

\section{Influencia de la Sentencia Penal sobre la Civil}

Lo que dice el Código Civil y Comercial es que la \textbf{inexistencia del hecho, la autoría y el no dolo penal} de responsabilidad determinados en la sentencia penal no pueden desconocerse en la sentencia civil. Es decir:

\begin{important}
Si la sentencia penal dice que el sindicado como responsable penal no participó, que no es autor del hecho, o que el hecho no ocurrió, esto no podrá ser desconocido en la sentencia civil. Deberá fundarse en consecuencia.
\end{important}

\subsection{Hipótesis de Absolución Penal}

Es posible que una absolución penal por ausencia de delito —porque el accionar no está tipificado en el Código Penal— sin embargo dé lugar a la reparación civil. En esos casos no habrá ninguna influencia de las consecuencias penales sobre las civiles.

\section{Sentencia Penal Posterior a la Civil}

\subsection{Régimen General}

Según el artículo 1780 del Código Civil y Comercial, la sentencia penal posterior a la sentencia civil no produce ningún efecto sobre ella, excepto en el caso de revisión.

La revisión procede cuando así se solicite y cuando la sentencia civil asigna alcances de cosa juzgada a cuestiones resueltas por la sentencia penal y ésta es revisada respecto a esas cuestiones.

\subsection{Casos de Revisión}

\begin{example}{Absolución penal posterior por inexistencia del hecho}
Si se dictó una sentencia civil condenando al demandado sobre la base de un factor de atribución objetivo (artículo 1774, inciso c, Código Civil y Comercial), y con posterioridad en la acción penal el demandado es absuelto por inexistencia del hecho en que se funda esa condena civil, o por no ser su autor, en ese caso procede la revisión de la condena civil. Aunque la sentencia penal llegó con posterioridad, está diciendo que el demandado que fue condenado en sede civil no fue autor del hecho, de modo que de ninguna manera podría quedar condenado.
\end{example}

% ============================================================================
% APÉNDICE: CUADRO CORNELL
% ============================================================================

\appendix

\chapter{Cuadro Cornell: Síntesis y Repaso}

\section{Preguntas, Respuestas y Resumen}

\begin{table}[H]
\centering
\begin{tabularx}{\textwidth}{
    >{\raggedright\arraybackslash}p{0.28\textwidth}
    >{\raggedright\arraybackslash}X
}
\toprule
\textbf{Preguntas / palabras clave} &
\textbf{Notas / respuestas} \\
\midrule

\textbf{¿Qué son las resoluciones judiciales y cuál es su clasificación?} &
Son las diversas formas en que se pronuncia el tribunal. Se clasifican en: providencias simples (mero trámite, 3 días, pueden firmarlas secretario, pro-secretario o jefe de despacho); resoluciones interlocutorias (deciden cuestiones incidentales con motivación, 10 días, solo el juez); sentencias homologatorias (modos anormales de terminación: allanamiento, desistimiento, conciliación); y sentencias definitivas de primera y segunda instancia (acto más complejo, art. 163 CPCC, estructura: resultando / considerando / fallo). \\

\textbf{¿Cómo se diferencia una providencia simple de una resolución interlocutoria?} &
Por la sustanciación, entendida como ``sustancia'' (contenido y motivación), no como ``traslado''. La providencia simple solo impulsa el trámite sin motivación; la interlocutoria requiere que el juez funde su decisión porque está resolviendo algo (admisión o rechazo de una excepción, medida cautelar, etc.). El ejemplo de la medida cautelar confirma que la sustanciación no equivale a traslado. \\

\textbf{¿Por qué es importante saber ante qué tipo de resolución se está?} &
Porque cada resolución judicial tiene determinados remedios o recursos disponibles en plazos específicos. No puede ``matarse un mosquito con un cañón'': si ante una providencia simple errónea se interpone el recurso equivocado, se pierde la posibilidad de corrección. \\

\textbf{¿Qué establece el art. 163 CPCC sobre la sentencia definitiva?} &
Establece 9 incisos: estructura (resultando, considerando, fallo); mención de las partes; fundamentos; posibilidad de usar presunciones (prueba indiciaria) basadas en hechos reales y concordantes bajo la regla de la sana crítica; congruencia (decisión expresa, positiva y precisa conforme a las pretensiones deducidas); hechos sobrevinientes debidamente probados; plazo para cumplir la condena; costas y honorarios. \\

\textbf{¿Qué es el principio de congruencia?} &
El juez debe hacerse cargo de las cuestiones introducidas en el proceso, pero no puede dar más de lo pedido ni algo distinto. Superar ese límite viola la congruencia y el derecho de defensa del demandado que no pudo contradecir lo que no se le reclamaba. \\

\textbf{¿Qué son los hechos sobrevinientes?} &
Son los hechos que suceden durante el transcurso del proceso y que modifican, extinguen o alteran la relación jurídica objeto del litigio. El juez debe adecuar la sentencia a esa nueva realidad; de lo contrario, la decisión quedaría discordante con la práctica. \\

\textbf{¿Cómo funciona la sentencia definitiva de segunda instancia?} &
Es elaborada por un acuerdo de tres jueces: un vocal preopinante y dos vocales que adhieren o disienten. Puede salir dos a uno. La redacción es en primera persona del singular (cada vocal emite su voto individual). No puede dictarse sin la presencia de los tres jueces. La cosa juzgada puede emanar de la sumatoria del fallo de primera y segunda instancia cuando la modificación fue parcial. \\

\textbf{¿Qué puede hacer el juez una vez dictada la sentencia (art. 166 CPCC)?} &
Solo puede enmendar errores puramente numéricos o errores materiales evidentes, y aclarar conceptos oscuros, sin alterar lo sustancial. Esto se hace de oficio (dentro de los 3 días) o a pedido de parte (aclaratoria: 3 días en 1ª instancia, 5 en 2ª). \\

\textbf{¿Qué es la cosa juzgada?} &
Es la autoridad que adquiere una decisión judicial cuando contra ella no hay más recursos disponibles o el plazo para interponerlos ha vencido sin que las partes los dedujesen. Produce dos efectos: (1) ingresa al patrimonio del vencedor el derecho a exigir el cumplimiento; (2) la realidad previa deja de existir y solo queda lo que emana de la sentencia. Los acuerdos transaccionales homologados tienen también efectos de cosa juzgada. \\

\textbf{¿Cuál es la diferencia entre cosa juzgada formal y material?} &
La cosa juzgada material impide reabrir cualquier cuestión del fondo del asunto en cualquier proceso posterior. La cosa juzgada formal solo impide la discusión en el mismo proceso, pero permite iniciar un juicio ordinario posterior sobre cuestiones que no pudieron debatirse en el proceso anterior. \\

\textbf{¿Cuáles son los límites subjetivos y objetivos de la cosa juzgada?} &
Subjetivos: la cosa juzgada solo alcanza a quienes fueron partes del proceso (actor, demandado, terceros intervinientes). No puede afectar a quienes no tuvieron la posibilidad de defenderse.

Objetivos: no solo las cuestiones debatidas, sino también las conexas, accesorias, subsidiarias o por continencia que pudieron plantearse y no se plantearon (art. 347, inc. 6º). La cosa juzgada puede ser ``implícita'' respecto de rubros accesorios no reclamados. \\

\textbf{¿En qué consiste la prelación de la acción penal sobre la civil?} &
Según el art. 1774 CCyC, si la acción penal precede o se intenta durante el curso de la civil, la sentencia civil debe suspenderse hasta la conclusión del proceso penal, salvo: (a) extinción de la acción penal; (b) frustración efectiva del derecho a ser indemnizado por dilación del proceso penal; (c) factor objetivo de responsabilidad (innovación del CCyC). \\

\textbf{¿Qué efectos produce la sentencia penal sobre la civil?} &
La sentencia penal anterior vincula a la civil: si determina la inexistencia del hecho, la falta de autoría o la ausencia de dolo, esas conclusiones no pueden desconocerse en la sentencia civil. La sentencia penal posterior no afecta a la civil ya dictada, salvo caso de revisión. \\

\midrule

\multicolumn{2}{p{\textwidth}}{
\textbf{Resumen:} Las resoluciones judiciales constituyen el lenguaje formal a través del cual el tribunal se comunica con las partes y avanza el proceso. Su clasificación —providencias simples, resoluciones interlocutorias, sentencias homologatorias y sentencias definitivas— determina qué recursos pueden interponerse, en qué plazo y ante quién. El deber de motivación garantiza el control democrático de los actos del Poder Judicial. La sentencia definitiva, regulada en el artículo 163 CPCC, debe ser expresa, positiva, precisa y congruente con las pretensiones deducidas, pudiendo el juez valerse de presunciones indiciarias para forjar su convicción. Una vez firme, la sentencia adquiere autoridad de cosa juzgada, que puede ser material (definitiva) o formal (limitada a lo discutible en ese proceso), con límites objetivos y subjetivos que alcanzan cuestiones conexas, accesorias y subsidiarias (art. 347, inc. 6º). La coexistencia del proceso civil y penal exige una coordinación normativa que, con el nuevo Código Civil y Comercial, admite la continuación de la acción civil frente a factores objetivos de responsabilidad sin necesidad de esperar la conclusión del proceso penal, evitando siempre el ``escándalo lógico'' de fallos contradictorios.
} \\

\bottomrule
\end{tabularx}
\end{table}

% ============================================================================
% FIN DEL DOCUMENTO
% ============================================================================

\end{document}
